(1) 实验室系中静质量为 $m_1$ 的物体以 $0<v<c$ 沿 $+x$ 运动，质量 $m_2$ 的另一物体最初静止，两者完全粘连。设相对论因子 $\gamma=(1-v^2/c^2)^{-1/2}$。由四动量守恒求合体速度 $V$。
(2) 由四动量范数推导合体的静质量 $M$，证明当 $v>0$ 时 $M>m_1+m_2$。
(3) 求碰撞转化成合体内部能的静能增加量，并用实验室能量、动量核对表达式。
(4) 展开 $v/c\ll1$ 的合体速度与新增内部能领先项，解释为什么不能以静质量简单相加替代相对论守恒。
